\section{Conclusion}\label{sec:conclusion}

We have presented \emph{Pratyabhij\~n\=a $\times$ Active Inference} as a falsifiable engineering hypothesis: that a recursive self-reflexivity layer, formalised as a Bayesian-model-reduction-gated re-entry into the cascade, produces measurable lifts on benchmarks that reward interpretive multiplicity. The plugin (\texttt{pratyabhijna-creative-engine}) is shipped with all artefacts needed to replay the experiments end-to-end.

The negative-result obligation in our pre-registration means we report what the data show whether or not it flatters the theory. The headline results, the within-PCE internal-validity test, and the sensitivity contrast against the same-substrate baseline together make a coherent story (regardless of sign) about what an explicit Pratyabhij\~n\=a-style cascade contributes to a small-LLM substrate.

\paragraph{What follows.} We call out three directions that fall out of the design but are out of scope here: (i) replacing \texttt{kriya} with a Claude/Sonnet-class polish model so the substrate is no longer the bottleneck on fluency; (ii) replacing the Hopfield store with a larger episodic memory and re-running consolidation cycles across multi-day sessions; (iii) building a continuously-running PCE agent that uses \texttt{vimarsa} as a controller for tool invocation rather than only for surface revision.

\paragraph{Acknowledgements.} The Pratyaksha context-engineering harness~\citep{PratyakshaContextEng2025} provided the ``witness invariant'' discipline that this whole project is held to. The \texttt{ralph-loop} pattern from Anthropic's Claude Code plugin gallery was adopted as the completion-promise spine. The \texttt{attractor-flow} plugin~\citep{AttractorFlow2025} provided the engineering template for ``MCP $\to$ real computation $\to$ audit-traceable output''.
